% Technische Systemdokumentation auf Basis des Springer-LLNCS-Templates.
% Projektstand: 02.09.2026
% Evaluationskapitel ergänzt: 12.09.2026; Versuchsplan, noch ohne Messungen.
% Vor der Abgabe bitte die mit TODO markierten Metadaten ersetzen.
\documentclass[runningheads]{llncs}

\usepackage[T1]{fontenc}
\usepackage[utf8]{inputenc}
\usepackage[ngerman]{babel}
\usepackage{graphicx}
\usepackage{amsmath}
\usepackage{amssymb}
\usepackage{booktabs}
\usepackage{tabularx}
\usepackage{xcolor}
\usepackage{url}
\usepackage{microtype}

\newcommand{\systemname}{Local Person Re-Identification MVP}
\newcommand{\todoentry}[1]{\textcolor{red!70!black}{[#1]}}
\newcommand{\yes}{\textbf{umgesetzt}}
\newcommand{\partlyimplemented}{\textbf{teilweise}}
\newcommand{\planned}{\textbf{geplant}}
\newcommand{\evalpending}{\textit{n.\,e.}}

\begin{document}

\title{Konzeption und Implementierung eines lokalen Systems zur Person-Re-Identifikation in Videodaten}
\titlerunning{Lokale Person-Re-Identifikation in Videodaten}

% TODO: Namen, ORCID und Reihenfolge anpassen.
\author{\todoentry{Vorname Nachname}}
\authorrunning{\todoentry{Nachname}}

% TODO: Hochschule, Fakultät, Studiengang, Modul und E-Mail einsetzen.
\institute{\todoentry{Hochschule für angewandte Wissenschaften}\\
\todoentry{Fakultät / Masterstudiengang / Modul}\\
\email{\todoentry{hochschul-email@example.org}}}

\maketitle

\begin{abstract}
Diese Arbeit dokumentiert die Konzeption und prototypische Implementierung eines lokal ausführbaren Systems zur Person-Re-Identifikation in Video- und Webcam-Daten. Ziel ist eine nachvollziehbare Ende-zu-Ende-Pipeline, die Personen detektiert, innerhalb eines Videostroms verfolgt, visuelle Merkmale als normalisierte Embeddings extrahiert und Beobachtungen unter synthetischen Kennungen persistiert. Das System kombiniert eine YOLO-basierte Personendetektion mit ByteTrack oder BoT-SORT, einen austauschbaren ReID-Encoder auf Basis von OSNet beziehungsweise HSV-Farbhistogrammen sowie eine lokale SQLite-Datenhaltung. Zur Verringerung fehlerhafter Identitätsprofile werden Bildausschnitte anhand von Schärfe, Helligkeit, Größe, Seitenverhältnis, Randkontakt und Detektionskonfidenz bewertet. Erst mehrere hinreichend gute Beobachtungen führen zu einem initialen Abgleich; spätere Profilaktualisierungen werden qualitätsgewichtet. Eine ergänzende Bewegungsanalyse berechnet Richtung, Geschwindigkeit im Bildraum und auffällige Positionssprünge, beeinflusst die Identitätsentscheidung im aktuellen Stand jedoch nicht. Eine Streamlit-Oberfläche ermöglicht Konfiguration, Live-Vorschau und Ergebnisinspektion. Die Analyse des Implementierungsstands zeigt einen funktionsfähigen Forschungsprototyp, jedoch noch keine produktionsreife oder empirisch validierte biometrische Anwendung. Insbesondere fehlen ein reproduzierbarer Benchmark, Zugriffskontrolle, Verschlüsselung, automatisierte Löschfristen und die vollständige Integration des vorbereiteten Fußballanalysemodus. Abschließend werden ein Evaluationsdesign, Datenschutzanforderungen und priorisierte Weiterentwicklungen abgeleitet.

\keywords{Person Re-Identification \and Computer Vision \and Multi-Object Tracking \and OSNet \and Applied Sciences}
\end{abstract}

\section{Einleitung}

Die Wiedererkennung von Personen über zeitlich getrennte Beobachtungen ist eine zentrale Aufgabe der angewandten Bildverarbeitung. Im Gegensatz zur reinen Objektdetektion, die Position und Klasse einer Person in einem Einzelbild bestimmt, soll Person Re-Identification (ReID) entscheiden, ob zwei visuelle Beobachtungen derselben Person zugeordnet werden können. Innerhalb eines kontinuierlichen Videos übernimmt zunächst ein Multi-Object-Tracker die kurzfristige Zuordnung. ReID wird relevant, sobald eine Spur abbricht, eine Person nach einer Verdeckung erneut erscheint oder Beobachtungen langfristig zusammengeführt werden sollen.

Das in dieser Arbeit untersuchte System \systemname{} ist als kleiner, lokal betreibbarer Forschungsprototyp ausgelegt. Es verwendet keine Klarnamen, sondern Kennungen wie \path{person_000001}. Die zentrale technische Fragestellung lautet:

\begin{quote}
Wie kann eine modular erweiterbare, lokal ausführbare Videoanalyse-Pipeline gestaltet werden, die Tracking und Person-Re-Identifikation verbindet und zugleich die Qualität der gespeicherten Identitätsrepräsentationen kontrolliert?
\end{quote}

Die Arbeit leistet vier Beiträge: Erstens wird die implementierte Architektur vom Videoeingang bis zur lokalen Persistenz beschrieben. Zweitens werden die tatsächlich eingesetzten Matching-, Qualitäts- und Bewegungsverfahren formalisiert. Drittens werden Implementierungsstand und technische Grenzen transparent bewertet. Viertens wird ein reproduzierbares Evaluationskonzept formuliert, ohne nicht erhobene Messwerte als Ergebnisse auszugeben.

Der dokumentierte Systemstand entspricht dem Repository in Version 0.1.0 vom 2.~September~2026. Aussagen über Funktionen beziehen sich auf den Quellcode dieses Stands. Das Evaluationskapitel wurde am 12.~September~2026 um ein szenariobasiertes Versuchsprotokoll ergänzt. Eine quantitative Gütebewertung mit annotiertem Testdatensatz liegt noch nicht vor.

\section{Zielsystem und Anforderungen}

\subsection{Systemgrenze}

Das System verarbeitet eine lokale Videodatei oder eine über OpenCV erreichbare Webcam. Es erzeugt ein annotiertes Ausgabevideo, speichert ausgewählte Personenausschnitte und führt synthetische Personenprofile sowie Ereignisse in SQLite. Die Modelle werden lokal ausgeführt. Bibliotheken oder Modellgewichte können bei Installation beziehungsweise Erstinbetriebnahme dennoch Netzwerkzugriffe benötigen; eine vollständig netzwerkisolierte Ausführung setzt daher vorab bereitgestellte Abhängigkeiten und Gewichte voraus.

Nicht zum aktuellen Funktionsumfang gehören eine Feststellung bürgerlicher Identitäten, Gesichtserkennung, ein Mehrkamera-Abgleich mit Zeitsynchronisation, eine Benutzer- und Rechteverwaltung oder die automatisierte Ableitung rechtlich relevanter Entscheidungen. Ebenso ist die vorbereitete Fußballanalyse noch nicht mit dem produktiven Ausführungspfad verdrahtet.

\subsection{Funktionale Anforderungen}

Aus Projektziel und Implementierung ergeben sich die in Tabelle~\ref{tab:functional-requirements} zusammengefassten Anforderungen.

\begin{table}[t]
\caption{Funktionale Anforderungen und Implementierungsstand.}
\label{tab:functional-requirements}
\centering
\small
\begin{tabularx}{\columnwidth}{@{}p{0.12\columnwidth}X p{0.20\columnwidth}@{}}
\toprule
ID & Anforderung & Stand \\
\midrule
FA-01 & Video-Upload und lokale Webcam als Eingabe & \yes \\
FA-02 & Erkennung und kurzfristiges Tracking von Personen & \yes \\
FA-03 & Erzeugung austauschbarer ReID-Embeddings & \yes \\
FA-04 & Ähnlichkeitssuche und Vergabe synthetischer IDs & \yes \\
FA-05 & Qualitätskontrollierte Profilaktualisierung & \yes \\
FA-06 & Live-Vorschau und annotiertes Ausgabevideo & \yes \\
FA-07 & Speicherung von Läufen, Personen und Ereignissen & \yes \\
FA-08 & Frei konfigurierbare Modus-Presets & \yes \\
FA-09 & Fußballereignisse und belastbare Spielstatistiken & \planned \\
\bottomrule
\end{tabularx}
\end{table}

\subsection{Nichtfunktionale Anforderungen}

Der Prototyp priorisiert lokale Ausführung, geringe Einstiegshürden und Erklärbarkeit. Zentrale Parameter sind über Benutzeroberfläche oder Kommandozeile konfigurierbar. Einzelne Komponenten sollen austauschbar bleiben, insbesondere Detektionsmodell, Tracker und ReID-Encoder. Fehlerhafte oder zu kleine Bildausschnitte dürfen gespeicherte Profile nicht ungeprüft verschlechtern. Für einen produktiven Betrieb wären darüber hinaus messbare Anforderungen an Latenz, Verfügbarkeit, Skalierbarkeit, Zugriffssicherheit, Auditierbarkeit und Löschfristen festzulegen; sie sind im MVP noch nicht erfüllt.

\section{Fachliche und wissenschaftliche Grundlagen}

\subsection{Tracking-by-Detection}

Die Pipeline folgt dem Tracking-by-Detection-Paradigma. Für jedes Frame liefert der Detektor Bounding Boxes; der Tracker verknüpft diese zu zeitlichen Spuren mit lokalen \emph{track IDs}. ByteTrack assoziiert dabei auch Detektionen niedriger Konfidenz, um unter anderem unterbrochene Trajektorien bei Verdeckungen zu reduzieren~\cite{zhang2022bytetrack}. BoT-SORT ergänzt robuste Bewegungsassoziation um Kamerabewegungskompensation und optional Erscheinungsmerkmale~\cite{aharon2022botsort}. Beide Tracker werden über den Tracking-Modus der Ultralytics-Bibliothek angesprochen~\cite{ultralytics-track}.

Eine Tracker-ID ist keine dauerhafte Personenidentität. Sie gilt primär innerhalb einer zusammenhängenden Sequenz und kann nach Verdeckungen, Bildschnitten oder Verlassen des Sichtfeldes wechseln. Das untersuchte System hält deshalb eine separate Abbildung von Tracker-IDs auf synthetische Personen-IDs.

\subsection{Person-Re-Identifikation}

ReID-Modelle bilden einen Personenausschnitt auf einen Merkmalsvektor ab. OSNet lernt Merkmale über mehrere räumliche Skalen und wurde speziell für Person-ReID entworfen~\cite{zhou2019osnet}. Der Adapter des Systems verwendet \path{osnet_x1_0} und erzeugt einen 512-dimensionalen Vektor. Als leichtgewichtige Demonstrationsalternative steht ein 32-dimensionaler HSV-Histogrammvektor mit 16 Bins für den Farbton sowie je acht Bins für Sättigung und Helligkeit zur Verfügung. Diese Alternative validiert den technischen Datenfluss, ist jedoch bei ähnlicher Kleidung, Beleuchtungsänderungen und Kleidungswechseln nicht belastbar.

Für ein Roh-Embedding $x \in \mathbb{R}^{d}$ wird die $L_2$-Normalisierung
\begin{equation}
z = \frac{x}{\lVert x \rVert_2}
\label{eq:normalization}
\end{equation}
verwendet. Die Ähnlichkeit zweier normalisierter Vektoren $z_q$ und $z_p$ ist damit das Skalarprodukt
\begin{equation}
s(z_q,z_p) = \frac{z_q^\top z_p}{\lVert z_q\rVert_2\lVert z_p\rVert_2}=z_q^\top z_p.
\label{eq:cosine}
\end{equation}
Ein Profil gilt als Treffer, wenn der größte zulässige Score den konfigurierten Schwellwert $\tau$ erreicht. Im Default-Modus ist $\tau=0{,}82$ gesetzt. Der Wert ist eine Konfiguration, kein empirisch nachgewiesenes Optimum.

\section{Systemarchitektur}

\subsection{Komponenten und Datenfluss}

Abbildung~\ref{fig:architecture} zeigt den realisierten Hauptpfad. Die Orchestrierung erfolgt synchron und frameweise in einem einzelnen Prozess.

\begin{figure*}[t]
\centering
\fbox{\begin{minipage}{0.96\textwidth}
\centering
\small
\textbf{Video-Upload / lokale Webcam}
$\longrightarrow$
\textbf{OpenCV Frame Reader}
$\longrightarrow$
\textbf{YOLO Person Detection}\\[2mm]
$\longrightarrow$
\textbf{ByteTrack / BoT-SORT}
$\longrightarrow$
\textbf{Crop + Quality Gate}
$\longrightarrow$
\textbf{OSNet / HSV-Encoder}\\[2mm]
$\longrightarrow$
\textbf{Cosine Matching}
$\longleftrightarrow$
\textbf{SQLite Profile und Events}
$\longrightarrow$
\textbf{Live- und MP4-Ausgabe}
\end{minipage}}
\caption{Logischer Datenfluss der implementierten ReID-Pipeline.}
\label{fig:architecture}
\end{figure*}

Die zentralen Verantwortlichkeiten sind wie folgt verteilt:

\begin{itemize}
\item \path{app/ui/streamlit_app.py}: interaktive Konfiguration, Video-Upload, Kamerasuche, Fortschritt, Vorschau und Tabellenansichten;
\item \path{app/main.py}: nicht-interaktiver Kommandozeileneinstieg;
\item \path{app/pipeline/detector_tracker.py}: Adapter für YOLO und den gewählten Tracker;
\item \path{app/pipeline/reid_encoder.py}: einheitliche Schnittstelle für OSNet und Farbhistogramme;
\item \path{app/pipeline/orchestrator.py}: Ablaufsteuerung, Qualitätsfilter, Matching, Profilaktualisierung und Ausgabe;
\item \path{app/storage/vector_store.py}: Schema, Ähnlichkeitssuche und SQLite-Persistenz;
\item \path{app/utils/}: Bildverarbeitung, Bewegung, Kamerazugriff und ID-Erzeugung;
\item \path{app/modes/}: integrierte und benutzerdefinierte Konfigurations-Presets.
\end{itemize}

\subsection{Ablauf einer Beobachtung}

Nach dem Öffnen der Quelle werden Bildrate und Auflösung gelesen; bei ungültiger Bildrate wird mit 25~FPS weitergerechnet. YOLO verarbeitet ausschließlich die COCO-Klasse~0 (Person). Der Tracker wird mit persistierendem Zustand frameweise aufgerufen. Für jede Detektion wird entschieden, ob eine ReID-Auswertung erforderlich ist: bei einer noch nicht zugeordneten Spur oder in einem periodischen Intervall von standardmäßig fünf Frames.

Ein neuer Track wird nicht anhand eines einzelnen Frames registriert. Zunächst werden mindestens drei akzeptierte Kandidaten gesammelt. Anschließend entsteht aus ihnen ein kombiniertes Embedding. Die Datenbank liefert entweder ein vorhandenes Profil oberhalb des Schwellwerts oder das System vergibt eine neue synthetische ID. Bereits zugeordnete Tracks werden nicht bei jedem Update erneut global gematcht; gute Folgebeobachtungen aktualisieren das zugeordnete Profil. Eine globale ReID-Entscheidung erfolgt erneut, wenn der Tracker für die Person eine neue Track-ID liefert.

Innerhalb eines Frames verhindert eine Ausschlussmenge, dass dieselbe gespeicherte Personen-ID gleichzeitig zwei neuen Tracks zugeordnet wird. Dieser Mechanismus reduziert offensichtliche Doppelbelegungen, ersetzt jedoch keine globale zeitliche Optimierung.

\subsection{Modus-System}

Der integrierte Modus \path{default} parametrisiert die generische Person-ReID. Der Modus \path{football_team_analysis} setzt unter anderem BoT-SORT, einen leicht höheren Match-Schwellwert und Feature-Flags für Ball, Spielfeld, Teams und Statistiken. Die Flags dokumentieren die Zielarchitektur, aktivieren im aktuellen Ausführungspfad aber noch keine vollständige Fußballanalyse. Benutzerdefinierte Presets werden als JSON unter \path{data/modes/custom_modes.json} gespeichert und dürfen integrierte Modi nicht überschreiben.

\section{Methodischer Entwurf}

\subsection{Qualitätsbewertung der Personenausschnitte}

Schlechte Ausschnitte können ein ReID-Profil dauerhaft verschieben. Deshalb berechnet das System für jeden geometrisch gültigen Crop sechs auf $[0,1]$ skalierte Teilwerte: Laplace-basierte Schärfe $q_b$, ausgewogene Helligkeit $q_h$, relative Crop-Größe $q_s$, körperähnliches Seitenverhältnis $q_a$, fehlenden Randkontakt $q_e$ und Detektionskonfidenz $q_d$. Der implementierte Gesamtscore lautet
\begin{equation}
Q = 0{,}25q_b + 0{,}15q_h + 0{,}25q_s + 0{,}15q_a + 0{,}10q_e + 0{,}10q_d.
\label{eq:quality}
\end{equation}
Crops unter $Q_{\mathrm{embed}}=0{,}55$ werden nicht encodiert. Für die Aktualisierung eines bestehenden Profils gilt der strengere Standardwert $Q_{\mathrm{update}}=0{,}65$. Die Gewichte und Grenzwerte sind nachvollziehbare Heuristiken des MVP und müssen anhand annotierter Daten kalibriert werden.

\subsection{Multi-Frame-Fusion und Profilaktualisierung}

Für $n$ akzeptierte Kandidaten mit normalisierten Embeddings $z_i$ und Qualitätswerten $Q_i$ wird das initiale Track-Embedding qualitätsgewichtet zusammengeführt:
\begin{equation}
\bar{z}=\operatorname{norm}\left(\frac{\sum_{i=1}^{n} Q_i z_i}{\sum_{i=1}^{n} Q_i}\right).
\label{eq:fusion}
\end{equation}
Als Snapshot wird der Kandidat mit dem höchsten Qualitätswert ausgewählt. Bei späteren Beobachtungen wird das gespeicherte Mittel $\mu$ mit bisheriger Gewichtssumme $W$ und neuem Vektor $z$ gewichtet fortgeschrieben:
\begin{equation}
\mu' = \operatorname{norm}\left(\frac{W\mu+wz}{W+w}\right), \qquad W'=W+w,
\label{eq:update}
\end{equation}
wobei $w=Q$ mit einem internen Minimalgewicht von 0,05 ist. Wird der Encoder gewechselt und stimmt die Dimension des neuen Vektors nicht mit dem gespeicherten Profil überein, ignoriert die Suche inkompatible Profile; bei einer expliziten Aktualisierung wird der alte Mittelwert ersetzt. Für wissenschaftliche Vergleichbarkeit wäre eine explizite Encoder-Version pro Profil robuster als dieses MVP-Verhalten.

\subsection{Bewegungsanalyse}

Die Mitte einer Bounding Box bildet die Position im Bildraum. Aus aufeinanderfolgenden Positionen werden Verschiebung, geglättete Richtung und Pixelgeschwindigkeit berechnet. Für die Richtungsdarstellung verwendet das System eine exponentielle Glättung
\begin{equation}
\hat{d}_t=\alpha d_t+(1-\alpha)\hat{d}_{t-1}
\label{eq:motion-smoothing}
\end{equation}
mit $\alpha=0{,}35$ im Default-Modus. Ein Sprung gilt als auffällig, wenn seine Distanz größer ist als
\begin{equation}
D_{\max}=\sqrt{W_f^2+H_f^2}\,\rho\sqrt{\Delta f},
\label{eq:jump}
\end{equation}
wobei $W_f$ und $H_f$ die Frame-Dimensionen, $\rho=0{,}20$ den konfigurierten Anteil der Bilddiagonale und $\Delta f$ den Frame-Abstand bezeichnen. Die Motion-Daten werden visualisiert und in Ereignis-Payloads gespeichert. Sie dienen gegenwärtig ausschließlich der Diagnose und sind keine harte Bedingung für das ReID-Matching.

\section{Implementierung}

\subsection{Technologiestack und Konfiguration}

Tabelle~\ref{tab:stack} beschreibt die wesentlichen Technologien. Die Abhängigkeiten sind mit Mindestversionen, aber ohne Lock-Datei festgelegt. Dadurch bleibt die Installation flexibel, ist jedoch nicht vollständig reproduzierbar.

\begin{table}[t]
\caption{Technologiestack des Prototyps.}
\label{tab:stack}
\centering
\small
\begin{tabularx}{\columnwidth}{@{}p{0.28\columnwidth}X@{}}
\toprule
Baustein & Implementierung \\
\midrule
Sprache & Python ab Version 3.10 \\
Benutzeroberfläche & Streamlit ab Version 1.34 \\
Videoverarbeitung & OpenCV ab Version 4.9 \\
Detektion/Tracking & Ultralytics ab Version 8.2 \\
ReID & Torchreid/OSNet, optional HSV-Histogramm \\
Numerik/Daten & NumPy und pandas \\
Persistenz & SQLite über Python-Standardbibliothek \\
Optionaler Dienst & Qdrant-Container, noch ohne Adapter \\
\bottomrule
\end{tabularx}
\end{table}

Die Standardparameter in Tabelle~\ref{tab:defaults} stammen unmittelbar aus der implementierten Default-Konfiguration. Die tatsächliche Laufkonfiguration wird zusammen mit dem Analyse-Lauf in der Datenbank gespeichert.

\begin{table}[t]
\caption{Zentrale Standardparameter des Default-Modus.}
\label{tab:defaults}
\centering
\small
\begin{tabularx}{\columnwidth}{@{}X p{0.26\columnwidth}@{}}
\toprule
Parameter & Standardwert \\
\midrule
Detektionsmodell & \path{yolov8n.pt} \\
Tracker & \path{bytetrack.yaml} \\
Encoder & \path{torchreid} / OSNet \\
Cosine-Schwellwert & 0,82 \\
Detektionskonfidenz & 0,35 \\
Eingangsgröße & 640 Pixel \\
ReID-Intervall & 5 Frames \\
Gute Initialframes & 3 \\
Minimale Crop-Größe & $30\times80$ Pixel \\
Maximale Frames & 500 \\
Live-Vorschau & jedes 10. Frame \\
\bottomrule
\end{tabularx}
\end{table}

\subsection{Persistenzmodell}

SQLite wird in diesem MVP zugleich als Metadaten- und einfacher Vektorspeicher verwendet. Die Ähnlichkeitssuche lädt alle Profile und berechnet die Cosine Similarity in Python. Dies vermeidet einen zusätzlichen Dienst, skaliert aber linear mit der Anzahl gespeicherter Personen. Tabelle~\ref{tab:data-model} zeigt die Tabellen und ihre Rolle.

\begin{table*}[t]
\caption{Logisches Datenmodell. Fußballtabellen sind vorbereitet, werden im Hauptpfad jedoch noch nicht befüllt.}
\label{tab:data-model}
\centering
\small
\begin{tabularx}{\textwidth}{@{}p{0.20\textwidth}p{0.16\textwidth}X p{0.13\textwidth}@{}}
\toprule
Tabelle & Schlüssel & Wesentliche Inhalte & Nutzung \\
\midrule
\path{persons} & \path{person_id} & Mittel-Embedding, Gewichtssumme, Beobachtungen, Zeitpunkte, bestes Snapshot & produktiv \\
\path{events} & \path{event_id} & Quelle, Frame, Track, Box, Match-Score, Qualitäts- und Bewegungsdaten & produktiv \\
\path{analysis_runs} & \path{run_id} & Modus, Quelle, FPS, Auflösung, Konfigurationsmetadaten & produktiv \\
\path{teams}, \path{players} & zusammengesetzt & Team- und Spielerstammdaten pro Lauf & vorbereitet \\
\path{player_frame_events} & \path{event_id} & Spielerposition, Geschwindigkeit, Distanz, Konfidenz & vorbereitet \\
\path{ball_frame_events} & \path{event_id} & Ballposition, Geschwindigkeit, nächster Spieler/Team & vorbereitet \\
\path{player_stats} & Lauf und Spieler & Sichtbarkeit, Distanz, Geschwindigkeit, Sprints, Ballnähe, Heatmap & vorbereitet \\
\bottomrule
\end{tabularx}
\end{table*}

Embeddings werden als JSON-kodierte Float-Arrays gespeichert. Akzeptierte Updates erzeugen Ereignisse und Snapshot-Dateien. Das Feld \path{best_snapshot_path} zeigt lediglich auf den qualitativ besten bekannten Ausschnitt; ältere Snapshot-Dateien werden dabei nicht automatisch gelöscht. Das Reset-Skript entfernt Datenbank, Snapshots und Ausgabevideos vollständig und legt die Verzeichnisse neu an.

\subsection{Benutzerschnittstellen und Betrieb}

Die Streamlit-Oberfläche bietet Moduswahl, Parameter, Upload, Kamerasuche, Testbild, Live-Anzeige, Fortschrittsinformation und tabellarische Einsicht in Läufe, Personen sowie Ereignisse. Da Streamlit als Client-Server-Anwendung arbeitet, greift die lokale Webcam auf dem Rechner des Streamlit-Prozesses zu und nicht automatisch auf eine Kamera des entfernten Browser-Clients~\cite{streamlit-architecture}.

Der interaktive Start erfolgt aus dem Projektverzeichnis mit
\begin{verbatim}
python -m streamlit run app/ui/streamlit_app.py
\end{verbatim}
Die Kommandozeile unterstützt zusätzlich Quelle, Modus, Modell, Tracker, Encoder, Schwellwert, maximale Framezahl und Gerät. Pfade sowie einige Defaultwerte lassen sich über Umgebungsvariablen konfigurieren. Der optionale Qdrant-Dienst in \path{docker-compose.yml} ist infrastrukturell vorbereitet, wird von der Anwendung aber nicht angesprochen.

\section{Evaluation}
\label{sec:evaluation}

Dieses Kapitel beschreibt das vorgesehene Verfahren zur empirischen Bewertung des Systems. Die Untersuchung verbindet kontrollierte Eigenaufnahmen mit einem externen Videodatensatz. Im Mittelpunkt stehen die Zuverlässigkeit der synthetischen Identitätszuordnung, ihre Robustheit gegenüber veränderten Aufnahmebedingungen und die technische Leistungsfähigkeit. Zum Zeitpunkt der Erstellung ist die Evaluation noch nicht durchgeführt. Die folgenden Angaben beschreiben das festzulegende Versuchsprotokoll; die Ergebnistabellen enthalten ausschließlich als solche gekennzeichnete Platzhalter.

\subsection{Untersuchungsziel und Forschungsfragen}

Die Evaluation betrachtet drei aufeinander aufbauende Verarbeitungsebenen: Personendetektion, kurzfristiges Tracking und persistente Re-Identifikation. Ein hoher Detektions-Recall belegt noch keine korrekte Identitätszuordnung. Ebenso kann ein Tracker eine Person zuverlässig verfolgen, während die Anwendung nach einem Wiedereintritt eine neue synthetische ID vergibt. Die getrennte Messung erlaubt es, Fehler ihrer jeweiligen Entstehungsstufe zuzuordnen.

Die Untersuchung wird durch vier Forschungsfragen strukturiert:
\begin{enumerate}
\item \textbf{F1 -- Erfassung und Kontinuität:} Wie zuverlässig werden eine beziehungsweise mehrere gleichzeitig sichtbare Personen detektiert und über die Zeit verfolgt?
\item \textbf{F2 -- Wiedererkennung:} Wie zuverlässig erhält eine bereits registrierte Person nach Verlassen des Sichtfelds oder nach einem neuen Programmlauf ihre bisherige synthetische ID? Wie häufig werden unbekannte Personen fälschlich bestehenden Profilen zugeordnet?
\item \textbf{F3 -- Robustheit und Komponentenbeitrag:} Wie wirken sich ähnliche Kleidung, Verdeckungen, Bewegung, Perspektive und Beleuchtung aus? Welchen messbaren Beitrag leisten Encoder, Quality Gate und Multi-Frame-Fusion?
\item \textbf{F4 -- Betriebseigenschaften:} Welche Verarbeitungsgeschwindigkeit, Verzögerung bis zur Identitätszuweisung und welchen Speicherbedarf erreicht das System auf der dokumentierten Hardware?
\end{enumerate}

Als gerichtete Arbeitshypothesen werden eine geringere Fehlzusammenführungsrate von OSNet gegenüber Farbhistogrammen bei ähnlicher Kleidung sowie weniger fehlerhafte Profilaktualisierungen durch die Qualitätsfilterung untersucht. Gleichzeitig kann die Qualitätsfilterung eine längere Wartezeit oder ausbleibende Zuweisungen verursachen. Beide Effekte werden gemeinsam ausgewertet. Diese Erwartungen sind keine vorweggenommenen Ergebnisse.

\subsection{Datengrundlage und Aufnahmebedingungen}

\subsubsection{Kontrollierte Eigenaufnahmen.}
Geplant ist eine Stichprobe von vier freiwillig teilnehmenden Personen: den Projektmitgliedern sowie gegebenenfalls ein bis zwei weiteren Personen mit dokumentiertem Einverständnis. Die Aufnahmen erfolgen in einem kontrollierten Innenraum. Referenzkennungen wie \path{GT_01} bis \path{GT_04} bezeichnen die tatsächlichen Personen innerhalb des Versuchs. Die Anwendung gibt ausschließlich ihre eigenen synthetischen IDs aus. Die Ground-Truth-Kennungen werden ihr bei der Verarbeitung nicht übergeben.

Als Ausgangsbedingung dienen eine fest aufgestellte Kamera, gleichbleibende Beleuchtung und gut unterscheidbare Kleidung. Kameraposition, Aufnahmeauflösung, Bildrate und markierte Laufwege werden dokumentiert. Zunächst werden kurze, separate Pilotaufnahmen erstellt. Anschließend sind pro festgelegter Testbedingung mindestens drei eigenständige Aufnahmen von etwa 45 bis 90 Sekunden vorgesehen; für den leeren Raum genügt dieselbe zuvor festgelegte Beobachtungsdauer ohne Personen. Bei mehreren Varianten eines Szenarios wird jede Variante wiederholt. Die endgültige Anzahl der Clips und Ereignisse wird im Versuchsverzeichnis festgehalten.

Die Wiederholungen variieren die Reihenfolge der Personen und die Rollen bei Kreuzungen oder Verdeckungen. Ein erneuter Programmdurchlauf desselben Videos zählt lediglich als technische Wiederholung und erhöht nicht die Zahl unabhängiger Aufnahmen. Vier Personen erlauben eine explorative Fallstudie unter kontrollierten Bedingungen, jedoch keine repräsentative Bewertung einer größeren Bevölkerung.

\subsubsection{Externe Videodaten.}
Als Ergänzung ist mindestens eine vorab ausgewählte, annotierte Sequenz aus MOT17 vorgesehen~\cite{dendorfer2021motchallenge,mot17-data}. Die Auswahl erfolgt nach Szeneneigenschaften, etwa höherer Personendichte oder Kamerabewegung, bevor die Systemergebnisse betrachtet werden. Verwendet werden Sequenzen mit zugänglichen Referenzannotationen aus dem veröffentlichten Trainingsbestand als externer Evaluationsbestand; auf diesen Sequenzen werden weder Modelle trainiert noch Parameter abgestimmt. Sequenzname, verwendeter Framebereich, Herkunft und Nutzungsbedingungen werden im Protokoll angegeben.

Die mehrfach angebotenen Detektorvarianten derselben MOT17-Sequenz stellen keine unabhängigen Videos dar. Da das System eigene YOLO-Detektionen erzeugt, wird jede ausgewählte Bildsequenz nur einmal gezählt. Die externe Auswertung erfolgt lokal und liefert eine zusätzliche Fallstudie; sie ist keine vollständige Benchmarkteilnahme. Externe Daten und Eigenaufnahmen werden getrennt berichtet. Ein frei abrufbares Video ohne Referenzannotation wäre zunächst nur für eine qualitative Demonstration geeignet.

\subsubsection{Umgang mit Aufnahmen und Kennungen.}
Vor der Aufnahme werden Zweck, Verarbeitung von Videos, Ausschnitten und Embeddings, Zugriffsberechtigte, Aufbewahrungsdauer sowie die geplante Verwendung in der Hochschularbeit mit den Teilnehmenden dokumentiert. Einwilligungsunterlagen werden getrennt von technischen Kennungen aufbewahrt. Die Veröffentlichung erkennbarer Beispielbilder wird gesondert geklärt. Synthetische IDs bewirken keine automatische Anonymisierung; die kontextbezogene Datenschutzprüfung wird in Abschnitt~\ref{sec:privacy} behandelt~\cite{edpb-video}.

\subsection{Szenarien und kontrollierte Variationen}

Die Tabellen~\ref{tab:eval-core-cases} und~\ref{tab:eval-extended-cases} operationalisieren die Use Cases. In den kontrollierten Tests wird gegenüber der Baseline möglichst jeweils ein Einflussfaktor verändert. Kombinationen, beispielsweise ähnliche Kleidung bei gleichzeitiger Verdeckung, werden als zusätzliche Belastungstests gesondert ausgewiesen.

\begin{table}[tbp]
\caption{Kernszenarien der kontrollierten Evaluation.}
\label{tab:eval-core-cases}
\centering
\small
\begin{tabularx}{\linewidth}{@{}p{0.10\linewidth}p{0.23\linewidth}X@{}}
\toprule
ID & Szenario & Vorgesehener Ablauf und Beobachtungsziel \\
\midrule
UC-01 & Einzelperson & Eine Person steht und geht bei freier Sicht. Ermittlung der Baseline für Erfassung, ID-Zuweisung und Kontinuität. \\
UC-02 & Mehrere Personen & Zwei beziehungsweise vier Personen erscheinen zunächst getrennt, anschließend gleichzeitig. Prüfung auf fehlende Personen und gemeinsame IDs. \\
UC-03 & Raum verlassen und zurückkehren & Nach einer Registrierungsphase verlassen Personen das Sichtfeld und kehren nach kurzen und längeren, protokollierten Abständen zurück. Prüfung der bisherigen Personen-ID. \\
UC-04 & Ähnliche Kleidung & Zwei Personen tragen ähnlich gefärbte Ober- und Unterteile. Vergleich mit unterscheidbarer Kleidung bei gleichen Laufwegen. \\
UC-05 & Kreuzung und Verdeckung & Zwei Personen kreuzen sich; zusätzlich verdeckt eine die andere kurz vollständig. Prüfung von Track-Abbrüchen und Identitätswechseln. \\
UC-06 & Bewegung & Stillstand, normales und schnelles Gehen sowie Richtungswechsel entlang markierter Wege. Prüfung der Kontinuität bei unterschiedlicher Bewegung. \\
\bottomrule
\end{tabularx}
\end{table}

\begin{table}[tbp]
\caption{Ergänzende Robustheits-, Negativ- und Persistenztests.}
\label{tab:eval-extended-cases}
\centering
\small
\begin{tabularx}{\linewidth}{@{}p{0.10\linewidth}p{0.23\linewidth}X@{}}
\toprule
ID & Szenario & Vorgesehener Ablauf und Beobachtungsziel \\
\midrule
UC-07 & Entfernung und Ansicht & Markierte Nah- und Fernpositionen sowie Vorder-, Seiten- und Rückenansicht. Varianten getrennt auswerten. \\
UC-08 & Beleuchtung & Gleicher Ablauf bei normalem Licht, reduziertem Licht und Gegenlicht. Prüfung von Erfassung und Wiedererkennung. \\
UC-09 & Teilverdeckung & Ober- oder Unterkörper wird durch Tisch oder Türrahmen verdeckt. Prüfung des Quality Gate und der ID-Abdeckung. \\
UC-10 & Unbekannte Person & Drei Personen werden registriert; die vierte erscheint erstmalig. Prüfung der Neuvergabe statt eines falschen Treffers. Rollen zwischen Aufnahmen wechseln. \\
UC-11 & Leerer Raum & Beobachtung ohne anwesende Personen. Zählung falscher Detektionen und angelegter Phantomprofile. \\
UC-12 & Neuer Programmlauf & Registrierung in Video A, neuer Prozess für Video B mit unveränderter Versuchsdatenbank. Prüfung persistenter Wiedererkennung. \\
UC-13 & Verändertes Erscheinungsbild & Registrierte Person kehrt mit abgelegter Jacke oder verändertem Oberteil zurück. Gesonderter Test der Systemgrenzen. \\
UC-14 & Externes Video & Vorab ausgewählte MOT17-Sequenz mit Referenzannotation. Prüfung der Übertragbarkeit auf eine andere Aufnahmeumgebung. \\
\bottomrule
\end{tabularx}
\end{table}

Für UC-03 werden kurze und lange Abwesenheiten relativ zur tatsächlich verwendeten Tracker-Konfiguration festgelegt. Mindestens eine Variante überschreitet dessen Aufbewahrungsdauer verlorener Tracks. Bleibt die ursprüngliche \path{track_id} erhalten, kann die Anwendung ihre bestehende Zuordnung weiterverwenden, ohne einen neuen ReID-Abgleich auszuführen. Daher werden Rückkehrereignisse mit beibehaltenem Track und Rückkehrereignisse mit neuem Track getrennt gezählt. UC-12 erzwingt zusätzlich einen frischen Tracker-Zustand. Nur dadurch lässt sich die Wirkung der persistenten ReID von der kurzfristigen Track-Fortsetzung unterscheiden.

\subsection{Referenzannotation und erforderliche Messdaten}

\subsubsection{Ground Truth.}
Für die ausgewählten Eigenaufnahmen werden alle auszuwertenden Frames mit Personenkennung, Bounding Box und Sichtbarkeitsstatus annotiert. Die Box umfasst den geschätzten vollständigen Körper innerhalb der Bildgrenzen; teilweise verdeckte Personen werden bei nachvollziehbarer Identität weiter annotiert. Vollständig verdeckte oder außerhalb des Bildes befindliche Personen erhalten in diesem Zeitraum keine auszuwertende Detektionsbox, behalten jedoch ihre Referenzidentität bei der Rückkehr. Tatsächlich uneindeutige Intervalle werden vor Betrachtung der Systemausgabe als nicht auswertbar markiert. Deren Dauer und Ausschlussgrund werden berichtet. Fehlerhafte oder schwierige Systemergebnisse sind kein Ausschlussgrund.

Ein- und Austritte sowie Anfang und Ende von Verdeckungen werden zusätzlich als Ereignisse markiert. Kritische Übergänge und eine vorab bestimmte Stichprobe der übrigen Frames werden von einem zweiten Projektmitglied überprüft. Eine automatische Interpolation von Boxen wird insbesondere bei Kreuzungen und Richtungswechseln manuell kontrolliert. Externe Sequenzen behalten ihr eigenes Annotations- und Ignore-Protokoll; eine Vergleichbarkeit mit den Eigenaufnahmen wird nicht allein aus identischen Metriknamen abgeleitet.

\subsubsection{Vorhersagen und Ereignisse.}
Für eine reproduzierbare Evaluation ist ein maschinenlesbarer Export pro Frame erforderlich. Er muss mindestens Sequenz, Lauf, Frameindex, Videozeit, Bounding Box, Detektionskonfidenz, \path{track_id}, \path{person_id} beziehungsweise einen fehlenden Zuweisungsstatus enthalten. Zusätzlich werden tatsächliche ReID-Entscheidungen mit Entscheidungsframe, neuem oder bestehendem Profil, aktuellem Ähnlichkeitsscore und Qualitätswert protokolliert. Frames ohne Detektionen bleiben über ein vollständiges Frameverzeichnis erkennbar.

Die bestehende Ereignistabelle protokolliert ausgewählte Matches und Profilaktualisierungen. Sie bildet nicht sämtliche Frame-Vorhersagen ab. Auch die Zähler für neu angelegte Personen und Matches enthalten keine Prüfung gegen Ground Truth. Sie genügen deshalb nicht zur Berechnung von Precision, Recall, IDF1 oder HOTA. Der beschriebene Export und die Auswertung sind vor Durchführung noch zu ergänzen.

Für eine isolierte Detektorbewertung müssen die Personendetektionen vor der Track-Zuordnung exportiert werden. Werden ausschließlich die Ausgaben des Tracker-Adapters verwendet, sind Precision und Recall ausdrücklich als Werte der kombinierten Detektions- und Tracking-Ausgabe zu kennzeichnen. Für Zeitmessungen ist außerdem der tatsächliche Entscheidungszeitpunkt maßgeblich: Das gespeicherte initiale Ereignis kann auf das frühere Frame des besten gepufferten Ausschnitts verweisen. Ein später erneut angezeigter Match-Score ist ohne neuen Abgleich keine unabhängige Messung.

\subsection{Metriken und Auswertungsregeln}

\subsubsection{Detektion.}
Vorhersage und Referenzbox werden für die eigene Detektionsauswertung je Frame eindeutig zugeordnet. Zulässig sind Paare mit einer Intersection over Union von mindestens 0,5:
\begin{equation}
\operatorname{IoU}(B_p,B_g)=
\frac{|B_p\cap B_g|}{|B_p\cup B_g|}.
\label{eq:eval-iou}
\end{equation}
Unter den zulässigen Paaren wird zunächst die größtmögliche Anzahl von Zuordnungen und anschließend die größte Gesamtüberlappung bestimmt. Ein zugeordnetes Paar zählt als True Positive (TP), eine übrige Vorhersage als False Positive (FP) und eine übrige Referenzbox als False Negative (FN). Daraus folgen
\begin{equation}
P_D=\frac{TP}{TP+FP},\qquad
R_D=\frac{TP}{TP+FN},\qquad
F_{1,D}=\frac{2TP}{2TP+FP+FN}.
\label{eq:eval-detection}
\end{equation}
Ergänzend werden die mittlere IoU der korrekten Detektionen und die Zahl falscher Detektionen pro Videominute berichtet. Bei einem Nenner von null wird eine Kennzahl als nicht definiert ausgewiesen. Insbesondere besitzt der vollständig leere Raum keinen sinnvoll berechenbaren Personen-Recall.

\subsubsection{Tracking.}
Die zeitliche Zuordnung der \path{track_id} wird mit IDF1 und HOTA bewertet. IDF1 basiert auf einer globalen Zuordnung von vorhergesagten und tatsächlichen Trajektorien~\cite{ristani2016identity}:
\begin{equation}
\operatorname{IDF1}=
\frac{2\,IDTP}{2\,IDTP+IDFP+IDFN}.
\label{eq:eval-idf1}
\end{equation}
Dabei bezeichnen $IDTP$, $IDFP$ und $IDFN$ identitätsbezogene korrekte, falsche und fehlende Zuordnungen der Standardmetrik. HOTA verbindet Detektions- und Assoziationsqualität über mehrere Lokalisierungsschwellen; ergänzend werden DetA und AssA zur Fehleranalyse betrachtet~\cite{luiten2021hota}. Absolute ID-Switches (IDSW) und Fragmentierungen (Frag) machen Identitätswechsel beziehungsweise Unterbrechungen sichtbar. MOTA wird als ergänzende Vergleichsgröße geführt.

Die Berechnung erfolgt mit einer festgehaltenen Version von TrackEval und dessen Standardprotokoll~\cite{trackeval}. Die separat festgelegte IoU-Regel der eigenen Detektionsauswertung ersetzt nicht die Regeln dieser Metriken. Tracker-IDs werden pro Sequenz ausgewertet. Ihre numerischen Werte müssen nicht den Ground-Truth-Kennungen entsprechen. Gute Trackingwerte allein belegen noch keine korrekte persistente \path{person_id}; diese wird durch die folgenden Ereignismetriken bewertet.

\subsubsection{Wiedereintritt und Identitätsentscheidung.}
Für jedes annotierte Ein- oder Rückkehrereignis wird ein Auswertungsfenster von zwei Sekunden ab dem ersten wieder sichtbaren Frame vorgesehen. Die betroffene Person muss danach mindestens zwei Sekunden auswertbar sichtbar bleiben. Kürzere Auftritte werden als eigene Kurzzeitfälle ausgewiesen. Die Fensterlänge wird im Pilotversuch überprüft und vor der abschließenden Evaluation festgeschrieben. Innerhalb des Fensters zählt die erste ausgegebene Personen-ID; bleibt eine Zuweisung aus, liegt ein zeitlich fehlgeschlagener Versuch vor. Spätere Korrekturen werden zusätzlich dokumentiert, ersetzen aber nicht das Ergebnis der ersten Entscheidung.

Die Referenz für eine bisherige Personen-ID wird anhand der Registrierungsphase vor dem jeweiligen Austritt festgelegt und danach nicht rückwirkend optimiert. Eine erfolgreiche Registrierung erfordert eine eindeutige ID ohne Vermischung mit einer anderen Referenzperson. Bereits fehlerhafte Registrierungen werden gezählt. Sie dürfen nicht durch nachträgliche Umbenennung der Vorhersagen korrigiert werden.

Seien $N_{\mathrm{return}}$ alle protokollgemäßen Rückkehrereignisse, $N_{\mathrm{registered}}$ die darunter liegenden Ereignisse mit zuvor erfolgreicher Registrierung und $N_{\mathrm{correct}}$ die rechtzeitig korrekt wiedererkannten Ereignisse. Die Ende-zu-Ende-Erfolgsrate und die bedingte Wiedererkennungsrate lauten
\begin{equation}
R_{\mathrm{return}}=
\frac{N_{\mathrm{correct}}}{N_{\mathrm{return}}},
\qquad
R_{\mathrm{registered}}=
\frac{N_{\mathrm{correct}}}{N_{\mathrm{registered}}}.
\label{eq:eval-reentry}
\end{equation}
Die erste Kennzahl berücksichtigt auch eine fehlgeschlagene Erstregistrierung. Die zweite beschreibt die Wiedererkennung bei bereits gültigem Profil. Beide werden mit ihren absoluten Zählern und getrennt nach unverändertem Track, neuem Track und neuem Prozess berichtet.

Für die Ereignisse mit zuvor registrierter Person sowie die gezielt unbekannten Eintritte werden zusätzlich ReID-Precision und ReID-Recall definiert. $C$ zählt korrekte Zuweisungen einer bestehenden ID an bekannte Personen, $W$ falsche Zuweisungen bestehender IDs, einschließlich Verwechslungen bekannter Personen und falscher Treffer bei unbekannten Personen. $K$ bezeichnet alle Rückkehrereignisse zuvor korrekt registrierter Personen:
\begin{equation}
P_{\mathrm{ReID}}=\frac{C}{C+W},\qquad
R_{\mathrm{ReID}}=\frac{C}{K},\qquad
F_{1,\mathrm{ReID}}=
\frac{2C}{C+W+K}.
\label{eq:eval-reid}
\end{equation}
Eine falsche bestehende ID bei einer bekannten Person beeinträchtigt somit Precision und Recall; eine neue oder fehlende ID beeinträchtigt ihren Recall. Die Zählform des F1 entspricht bei definierten positiven Werten ihrem harmonischen Mittel und ergibt bei vorhandenen Versuchen ohne korrekten Treffer null. Ein Nenner von null wird jeweils als nicht definiert ausgewiesen. Diese ereignisbezogenen Größen sind keine Ersatzbezeichnungen für die standardisierte Trackingmetrik IDF1.

Für unbekannte Eintritte wird der Anteil rechtzeitig vergebener neuer, nicht bereits belegter IDs an allen solchen Eintritten ausgewiesen. Zusätzlich wird $W/N_{\mathrm{events}}$ als ereignisbezogene Fehlzusammenführungsrate berichtet, wobei $N_{\mathrm{events}}$ die Summe der zuvor registrierten Rückkehrereignisse und unbekannten Eintritte ist. Die Aufspaltungsrate ergibt sich aus Rückkehrereignissen registrierter Personen mit einer neuen ID, geteilt durch $K$. Fehlende Zuweisungen werden gesondert gezählt. Bei UC-11 werden die erzeugten Phantomprofile absolut und pro Minute angegeben.

\subsubsection{Abdeckung und Verzögerung.}
Die ID-Abdeckung ist der Anteil auswertbarer Personen-Frames, in denen eine passende Box mit irgendeiner synthetischen Personen-ID vorliegt. Sie wird zusammen mit den Fehlzuordnungen berichtet, da eine hohe Abdeckung auch mit falschen IDs erreicht werden kann. Die Zeit bis zur Identifikation wird in Videozeit zwischen dem ersten sichtbaren Frame und der tatsächlichen ID-Entscheidung gemessen. Ausgewiesen werden Median und 95.~Perzentil für erfolgreiche Fälle sowie Anzahl und Anteil nicht rechtzeitig identifizierter Fälle. Dadurch bleiben Ablehnungen durch das Quality Gate sichtbar.

\subsubsection{Laufzeit und Speicher.}
Für $N$ tatsächlich verarbeitete Frames, die Verarbeitungsdauer $T_{\mathrm{proc}}$ und die ausgewertete Videodauer $T_{\mathrm{video}}$ werden
\begin{equation}
\operatorname{FPS}_{\mathrm{proc}}=\frac{N}{T_{\mathrm{proc}}},
\qquad
\operatorname{RTF}=\frac{T_{\mathrm{proc}}}{T_{\mathrm{video}}}
\label{eq:eval-runtime}
\end{equation}
berechnet. Ein Real-Time-Faktor unter eins bezeichnet eine Verarbeitung schneller als die Videodauer, belegt aber allein keine verzögerungsarme Live-Anzeige. Zusätzlich werden Median und 95.~Perzentil der Frame-Verarbeitungszeit, maximaler Arbeitsspeicher und gegebenenfalls GPU-Speicher erfasst. Das Wachstum von Datenbank und Snapshot-Verzeichnis wird als Byte-Differenz pro Videominute gemessen.

Modellladezeit und initiale Aufwärmphase werden separat protokolliert. Die gemessene Verarbeitung umfasst Lesen, Detektion, Tracking, ReID, Persistenz und Videoausgabe. Der Zustand der Live-Vorschau bleibt innerhalb eines Vergleichs konstant und wird angegeben. Hardware, GPU-Synchronisierung bei Zeitmessungen, Hintergrundlast und tatsächliche Framezahl werden dokumentiert; aus der nominellen Kamera-Bildrate wird keine Verarbeitungsleistung abgeleitet.

Rank-1 und mAP sind für ein ergänzendes Retrieval-Experiment mit fest definierter Query/Gallery-Aufteilung geeignet, beispielsweise auf Market-1501~\cite{zheng2015market1501}. Für den vorgesehenen Online-Versuch mit vier Personen stehen die oben definierten Identitätsereignisse im Vordergrund. Eine ROC-Kurve oder Equal Error Rate setzt zusätzlich ein festgelegtes Paarvergleichsprotokoll voraus und wird nicht aus wiederholt angezeigten Match-Scores abgeleitet.

\subsection{Versuchsablauf und Vergleichskonfigurationen}

\subsubsection{Kalibrierung und Testtrennung.}
Pilot- und Kalibrierungsaufnahmen werden vollständig von den abschließenden Testclips getrennt. Benachbarte Frames derselben Aufnahme werden nicht auf beide Mengen verteilt. Verwendet werden vortrainierte Modelle; ein zusätzliches Training ist in diesem Versuchsplan nicht vorgesehen. Da dieselben vier Personen in mehreren Aufnahmen auftreten können, prüft die Studie vor allem neue Situationen dieser Stichprobe. Eine identitätsunabhängige Generalisierung lässt sich daraus nicht ableiten.

Als Referenz dient die Standardkonfiguration aus Tabelle~\ref{tab:defaults}. Vor den Testläufen wird die Framebegrenzung so gesetzt, dass jedes ausgewählte Video vollständig verarbeitet wird; der Standardwert 500 würde längere Clips abschneiden. Modellgewichte, Bibliotheksversionen, Tracker-YAML, Geräteauswahl und alle wirksamen Konfigurationswerte werden archiviert. Ein eigener Datenspeicher pro Versuch verhindert eine Vermischung mit bisherigen Demo-Läufen.

Jeder unabhängige Konfigurationsvergleich beginnt mit leerer Versuchsdatenbank und neu initialisiertem Tracker. Bei Rückkehrversuchen bleibt die Datenbank innerhalb der zusammengehörigen Aufnahme erhalten. UC-12 verwendet ein zusammengehöriges Video-Paar: Beide Videos bilden eine Versuchseinheit, ausschließlich der Prozess wird dazwischen neu gestartet. Beim Encoderwechsel wird die Registrierung aus denselben Bildern mit dem jeweiligen Encoder neu aufgebaut; Profile unterschiedlicher Merkmalsräume werden nicht vermischt.

\begin{table}[tbp]
\caption{Geplante Ablationen. Jede Variante wird mit derselben Referenz verglichen.}
\label{tab:eval-ablations}
\centering
\small
\begin{tabularx}{\linewidth}{@{}p{0.10\linewidth}p{0.28\linewidth}X@{}}
\toprule
ID & Geänderte Einstellung & Aussageziel \\
\midrule
B0 & Standardkonfiguration & Referenz für alle Einzelvergleiche. \\
A1 & HSV statt OSNet & Beitrag des Encoders, insbesondere bei ähnlicher Kleidung. \\
A2 & BoT-SORT statt ByteTrack & Beitrag des Trackers bei Kreuzung und Verdeckung. \\
A3 & Beide Qualitätsschwellen auf null & Beitrag der qualitätsabhängigen Annahme und Ablehnung. Geometrische Mindestgrößen bleiben bestehen. \\
A4 & Ein statt drei Initialframes & Beitrag der Initialfusion und Auswirkung auf die Zuweisungszeit. \\
A5 & Schwellwertvariation & Verhältnis von falschen Treffern, Aufspaltungen und ausbleibenden Treffern. \\
\bottomrule
\end{tabularx}
\end{table}

Bei A3 bleiben qualitätsabhängige Gewichtung und Snapshot-Auswahl erhalten. Die Variante untersucht deshalb gezielt die Filtergrenzen und nicht den vollständigen Verzicht auf Qualitätsinformation. A4 verändert zugleich die Anzahl der Initialbeobachtungen und die frühestmögliche Entscheidung; beide Wirkungen sind bei der Interpretation zu berücksichtigen. Die gesamte Konfiguration eines anderen Betriebsmodus wird für A2 nicht übernommen, da dies mehrere Parameter gleichzeitig ändern würde.

Für A5 ist zunächst eine auf Kalibrierungsdaten ausgewertete Schwellwertmenge von 0,70; 0,75; 0,80; 0,82; 0,85; 0,90 und 0,95 vorgesehen. Pro Encoder wird der Wert mit dem höchsten ereignisbezogenen ReID-F1 gewählt; bei Gleichstand entscheidet die höhere Precision, danach der höhere Schwellwert. Die Zusammenstellung bekannter und unbekannter Eintritte bleibt konstant. Auf dem Testbestand werden die unveränderte Referenz sowie die vorab gewählte Konfiguration berichtet. Eine nachträgliche Wahl des besten Testwerts ist ausgeschlossen. Der Vergleich bei gemeinsamem Schwellwert und der Vergleich nach encoderbezogener Kalibrierung werden getrennt interpretiert.

\subsubsection{Ausführung und Aggregation.}
Für jede Aufnahme werden die festgelegten Konfigurationen in dokumentierter, wechselnder Reihenfolge ausgeführt. Mindestens drei technische Wiederholungen dienen der Laufzeitmessung. Für Qualitätskennzahlen wird die Aufnahme als Vergleichseinheit behandelt; identische Wiederholungen werden nicht zu zusätzlichen unabhängigen Testdaten erklärt.

Pro Use Case werden Einzelwerte je Clip sowie Mittelwert und Standardabweichung ausgewiesen. Ereignisraten werden zusätzlich als Zähler/Nenner angegeben; für eine gepoolte Rate werden zuerst die Ereigniszahlen addiert. Ein gleichgewichtiges Mittel der Clip-Raten wird ausdrücklich als solches bezeichnet. HOTA und IDF1 werden mit der Aggregationslogik des Evaluators zusammengeführt. Eigenaufnahmen und externer Bestand erhalten keine gemeinsame Gesamtpunktzahl.

Die Ablationen werden paarweise auf denselben Clips verglichen. Berichtenswert sind absolute Kennzahlen und deren Differenzen gegenüber B0, bei Prozentwerten in Prozentpunkten. Aufgrund der kleinen Stichprobe stehen deskriptive Aussagen im Vordergrund. Einzelne Frames und mehrere Ereignisse derselben Person sind statistisch abhängig; sie dürfen keine künstlich großen Stichproben für Signifikanztests bilden. Falls bei größerem Datenumfang Konfidenzintervalle berechnet werden, erfolgt das Resampling über unabhängige Aufnahmesitzungen und nicht über Einzelbilder.

\subsection{Darstellung und Interpretation der Ergebnisse}

Die Tabellen~\ref{tab:eval-result-tracking} bis~\ref{tab:eval-result-ablation} legen die Ergebnisdarstellung fest. \evalpending{} bedeutet \emph{noch nicht erhoben}; nicht anwendbare Größen werden später mit einem Gedankenstrich, rechnerisch undefinierte Werte mit \emph{n.\,d.} gekennzeichnet. Insbesondere steht kein Platzhalter für den Messwert null. Nach Durchführung werden die Tabellen bei mehreren Varianten um Bedingung, Clipzahl und Streuung erweitert. Numerische Akzeptanzgrenzen wären vor der Testauswertung anhand des Anwendungsziels festzulegen; aus noch nicht vorhandenen Daten wird keine bestandene Systemevaluation abgeleitet.

\begin{table}[tbp]
\caption{Vorlage: Detektions- und Trackingwerte der Referenzkonfiguration. $P_D$, $R_D$, IDF1 und HOTA in Prozent; IDSW und Frag als absolute Anzahlen.}
\label{tab:eval-result-tracking}
\centering
\small
\begin{tabularx}{\linewidth}{@{}Xrrrrrr@{}}
\toprule
Szenario & $P_D$ & $R_D$ & IDF1 & HOTA & IDSW & Frag \\
\midrule
UC-01 & \evalpending & \evalpending & \evalpending & \evalpending & \evalpending & \evalpending \\
UC-02 & \evalpending & \evalpending & \evalpending & \evalpending & \evalpending & \evalpending \\
UC-03 & \evalpending & \evalpending & \evalpending & \evalpending & \evalpending & \evalpending \\
UC-04 & \evalpending & \evalpending & \evalpending & \evalpending & \evalpending & \evalpending \\
UC-05 & \evalpending & \evalpending & \evalpending & \evalpending & \evalpending & \evalpending \\
UC-06 & \evalpending & \evalpending & \evalpending & \evalpending & \evalpending & \evalpending \\
UC-07 & \evalpending & \evalpending & \evalpending & \evalpending & \evalpending & \evalpending \\
UC-08 & \evalpending & \evalpending & \evalpending & \evalpending & \evalpending & \evalpending \\
UC-09 & \evalpending & \evalpending & \evalpending & \evalpending & \evalpending & \evalpending \\
UC-10 & \evalpending & \evalpending & \evalpending & \evalpending & \evalpending & \evalpending \\
UC-12 & \evalpending & \evalpending & \evalpending & \evalpending & \evalpending & \evalpending \\
UC-13 & \evalpending & \evalpending & \evalpending & \evalpending & \evalpending & \evalpending \\
UC-14 (extern) & \evalpending & \evalpending & \evalpending & \evalpending & \evalpending & \evalpending \\
\bottomrule
\end{tabularx}
\end{table}

\begin{table}[tbp]
\caption{Vorlage: Identitätsereignisse. $R$ bezeichnet Rückkehr, $U$ unbekannten Eintritt; Ergebnisse pro definierter Bedingung.}
\label{tab:eval-result-reid}
\centering
\small
\begin{tabularx}{\linewidth}{@{}Xrrrr@{}}
\toprule
Ereignismenge & Anzahl & \shortstack{Korrekt\\innerhalb 2 s} & \shortstack{Falsche\\bestehende ID} & \shortstack{Keine ID\\innerhalb 2 s} \\
\midrule
UC-03, $R$, alter Track & \evalpending & \evalpending & \evalpending & \evalpending \\
UC-03, $R$, neuer Track & \evalpending & \evalpending & \evalpending & \evalpending \\
UC-10, $U$ & \evalpending & \evalpending & \evalpending & \evalpending \\
UC-12, $R$, Neustart & \evalpending & \evalpending & \evalpending & \evalpending \\
UC-13, $R$, Kleidung & \evalpending & \evalpending & \evalpending & \evalpending \\
\bottomrule
\end{tabularx}
\end{table}

Für Rückkehrereignisse werden Aufspaltungen als weitere Ergebniskategorie ergänzt; bei zuvor fehlgeschlagener Registrierung wird der Ausgang gesondert ausgewiesen. Zusätzlich werden $N_{\mathrm{registered}}$, $R_{\mathrm{return}}$, $R_{\mathrm{registered}}$, ReID-Precision/-Recall/-F1 sowie ID-Abdeckung berichtet. Für UC-10 bedeutet eine korrekte Entscheidung eine neue ID, bei den Rückkehrereignissen dagegen die ursprüngliche ID. Für UC-12 werden Trackingwerte pro Einzelvideo und Wiedererkennungswerte über das Video-Paar angegeben.

\begin{table}[tbp]
\caption{Vorlage: technische Leistung und Negativkontrolle. Alle Messwerte stehen aus.}
\label{tab:eval-result-runtime}
\centering
\small
\begin{tabularx}{\linewidth}{@{}Xll@{}}
\toprule
Kenngröße & Einheit & Messwert \\
\midrule
Verarbeitungsrate & Frames/s & \evalpending \\
Real-Time-Faktor & dimensionslos & \evalpending \\
Frame-Verarbeitungszeit, Median / P95 & ms & \evalpending \\
Modellinitialisierung & s & \evalpending \\
Identifikationszeit, Median / P95 & s Videozeit & \evalpending \\
Nicht rechtzeitig identifizierte Ereignisse & Anzahl / Gesamt & \evalpending \\
Maximaler RAM / GPU-Speicher & MiB & \evalpending \\
Datenbank- / Snapshot-Wachstum & MiB/Videominute & \evalpending \\
UC-11: falsche Detektionen & Anzahl / Minute & \evalpending \\
UC-11: Phantomprofile & Anzahl / Minute & \evalpending \\
\bottomrule
\end{tabularx}
\end{table}

\begin{table}[tbp]
\caption{Vorlage: paarweise Komponentenvergleiche. Alle Änderungen beziehen sich auf B0 und denselben Clipbestand.}
\label{tab:eval-result-ablation}
\centering
\small
\begin{tabularx}{\linewidth}{@{}Xrrrr@{}}
\toprule
Variante & \shortstack{ReID-F1\\(\%)} & \shortstack{IDF1\\(\%)} & \shortstack{Fehlmerge\\(\%)} & \shortstack{FPS\\(Frames/s)} \\
\midrule
B0: Referenz & \evalpending & \evalpending & \evalpending & \evalpending \\
A1: Farbhistogramm & \evalpending & \evalpending & \evalpending & \evalpending \\
A2: BoT-SORT & \evalpending & \evalpending & \evalpending & \evalpending \\
A3: Filtergrenzen null & \evalpending & \evalpending & \evalpending & \evalpending \\
A4: Ein Initialframe & \evalpending & \evalpending & \evalpending & \evalpending \\
A5: Kalibrierter Wert & \evalpending & \evalpending & \evalpending & \evalpending \\
\bottomrule
\end{tabularx}
\end{table}

Die spätere Diskussion wird entlang der Forschungsfragen F1 bis F4 geführt. Zu F1 werden Detektionsfehler von Assoziationsfehlern getrennt. Zu F2 werden korrekte Rückkehr, falsche Zusammenführung, Aufspaltung und fehlende Entscheidung verglichen. Zu F3 werden Veränderungen gegenüber der Baseline und die paarweisen Ablationen interpretiert. Zu F4 wird geprüft, ob die erzielte Qualität mit der gemessenen Verarbeitungsgeschwindigkeit vereinbar ist. Insbesondere wird eine geringere Fehlerrate nicht als Verbesserung gewertet, ohne gleichzeitig eine mögliche Abnahme der ID-Abdeckung oder Zunahme der Verzögerung zu betrachten.

Die qualitative Fehleranalyse ergänzt die Kennzahlen durch dokumentierte Zeitpunkte und Ursachen: fehlende Detektion, unpassender Crop, Ablehnung durch das Quality Gate, Track-Wechsel oder falscher ReID-Treffer. Pro Fehlerkategorie werden Vorkommen gezählt und Beispiele nach einem vorab festgelegten Auswahlprinzip beschrieben. Solange keine Messungen vorliegen, bleiben sowohl die Bestätigung der Hypothesen als auch Aussagen zu Überlegenheit oder Echtzeitfähigkeit offen.

\subsection{Aussagegrenzen und technische Vorprüfung}

Die Übertragbarkeit ist durch die geringe Personenzahl, den kontrollierten Innenraum und möglicherweise ähnliche Aufnahmebedingungen begrenzt. Ein externes Video erweitert die beobachteten Bedingungen, ermöglicht aber noch keine allgemeine Aussage über andere Kameras, Kleidung oder Bevölkerungsgruppen. Für vortrainierte Modellgewichte wird deren Herkunft dokumentiert; eine unbekannte Überlappung ihrer Trainingsdaten mit dem externen Benchmark wird als Einschränkung benannt.

Weitere Risiken für die Validität sind Annotationfehler, eine unkontrolliert vorbelegte Datenbank, wechselnde Softwareversionen, nicht gespeicherte Konfigurationswerte und eine Vermischung von Videozeit und Rechenzeit. Die festgelegten Referenzregeln, Versuchsdatenbanken und getrennten Messgrößen adressieren diese Risiken. Aussagen über statistische Signifikanz oder repräsentative Fairness sind mit dem geplanten Umfang nicht vorgesehen.

Im Rahmen der ursprünglichen Dokumentation vom 2.~September~2026 wurden die Dateien unter \path{app} und \path{scripts} mit Python~3.13.1 syntaktisch kompiliert und Kernabhängigkeiten importiert. Dies ist eine technische Vorprüfung dieses damaligen Stands und kein Ergebnis der hier geplanten Evaluation. Vor der Datenauswertung müssen der Frame-Export, die Referenzannotation und die Metrikberechnung einschließlich einfacher kontrollierbarer Fehlerfälle überprüft werden. Erst anschließend können die ausstehenden Tabellen mit tatsächlich erhobenen Werten gefüllt werden.

\section{Datenschutz, Sicherheit und verantwortliche Nutzung}
\label{sec:privacy}

Synthetische IDs verhindern die Speicherung von Klarnamen, bewirken aber nicht automatisch Anonymität. Videoausschnitte, Bewegungsdaten und Embeddings können weiterhin einer identifizierbaren Person zugeordnet sein. Bei technischer Verarbeitung körperlicher Merkmale zum Zweck eindeutiger Identifizierung kann zudem die Einordnung als biometrische Daten nach Art.~4 Nr.~14 und Art.~9 DSGVO relevant werden~\cite{gdpr,edpb-video}. Ob und auf welcher Rechtsgrundlage ein konkreter Einsatz zulässig ist, hängt vom Nutzungskontext ab und muss vor einer Erhebung fachkundig geprüft werden.

Für Hochschulversuche sind mindestens folgende Maßnahmen vorzusehen:

\begin{itemize}
\item dokumentierter Zweck, Datenkategorien, Rechtsgrundlage und Verantwortlichkeit;
\item informierte Einwilligung, sofern sie die geeignete Rechtsgrundlage des Versuchs ist;
\item Datensparsamkeit, kurze Löschfristen und ein nachweisbarer Löschprozess;
\item getrennte Speicherung von Einwilligungsunterlagen und technischen Kennungen;
\item Zugriffsschutz, Verschlüsselung ruhender Daten und verschlüsselte Übertragung bei nicht rein lokalem Betrieb;
\item Protokollierung administrativer Zugriffe sowie Überprüfung, ob eine Datenschutz-Folgenabschätzung erforderlich ist;
\item keine Tests mit unbeteiligten Personen oder unkontrollierter öffentlicher Videoüberwachung.
\end{itemize}

Der aktuelle Prototyp erfüllt davon lediglich Teilaspekte: lokale Verarbeitung, synthetische IDs und einen vollständigen manuellen Reset. Es fehlen Authentifizierung, Rollen, Verschlüsselung, automatische Aufbewahrungsfristen, selektive Löschung, Audit-Logging und eine formale Einwilligungsverwaltung. Er darf deshalb nur in einer kontrollierten Forschungsumgebung mit berechtigtem Material eingesetzt werden.

Auch Fairness ist experimentell zu untersuchen. Die Qualität von Detektion und ReID kann mit Perspektive, Körpergröße im Bild, Verdeckung, Kleidung und Umgebungsbedingungen variieren. Gruppenbezogene Leistungsunterschiede dürfen mangels geeigneter, rechtmäßig erhobener Daten nicht angenommen oder ausgeschlossen werden. Ergebnisse des Prototyps sind nicht als alleinige Grundlage für Entscheidungen über Personen geeignet.

\section{Grenzen und Risiken des aktuellen Entwurfs}

\begin{itemize}
\item \textbf{Fehlende empirische Validierung:} Die konfigurierten Schwellwerte und Qualitätsgewichte sind Heuristiken.
\item \textbf{Domänenverschiebung:} Ein vortrainiertes OSNet kann in stark abweichenden Kameraszenen an Trennschärfe verlieren.
\item \textbf{Lineare Suche:} SQLite und Python-basierte Vektorvergleiche skalieren nicht für große Galerien.
\item \textbf{Single-User-Charakter:} ID-Erzeugung über die aktuelle Profilanzahl und fehlende Prozesskoordination sind nicht für parallele Läufe ausgelegt.
\item \textbf{Fehlende Profilversionierung:} Encodername, Dimension und Modellversion sind nicht direkt am Personenprofil gespeichert.
\item \textbf{Ungebremstes Datenwachstum:} Akzeptierte Folgebeobachtungen können zahlreiche Snapshot-Dateien und Eventzeilen erzeugen.
\item \textbf{Keine Transaktionsklammer für einen Lauf:} Ein abgebrochener Lauf kann bereits Metadaten und Teilergebnisse hinterlassen.
\item \textbf{Motion nur diagnostisch:} Große Sprünge werden markiert, korrigieren aber keine Identitätszuordnung.
\item \textbf{Videocodec und Hardware:} MP4-Ausgabe, Kamerabackend und Beschleunigung hängen von der lokalen Installation ab.
\item \textbf{Reproduzierbarkeit:} Mindestversionen und das Container-Tag \path{latest} erlauben unkontrollierte Versionsänderungen.
\end{itemize}

Diese Grenzen sind für einen explorativen Prototyp vertretbar, müssen aber vor Feldtests oder produktiver Nutzung adressiert werden.

\section{Erweiterung zur Fußballanalyse}

Das Repository enthält eine vorbereitete Domänenerweiterung. Ein einfacher Teamklassifikator bewertet die mittlere Farbe des Oberkörperbereichs, ein Pitch-Mapper kann aus mindestens vier Punktkorrespondenzen eine Homographie bestimmen, und ein Statistikakkumulator kann aus metrischen Spielfeldkoordinaten Sichtzeit, Distanz, Geschwindigkeit, Sprints und Heatmap-Punkte ableiten. Weiterhin existieren Datenmodelle für Ball- und Spielerereignisse.

Der eigentliche \path{FootballAnalysisPipeline}-Typ enthält jedoch noch keine zusätzliche Orchestrierung; die UI instanziiert unabhängig vom Modus die generische \path{PersonReIdPipeline}. Der Balldetektor liefert als Platzhalter stets eine leere Liste. Teamklassifikation, Homographie und Statistikakkumulation werden im Hauptpfad nicht aufgerufen. Der Modus ist somit eine Architekturvorbereitung und kein fertig implementiertes Fußballanalysesystem.

Die empfohlene Reihenfolge der Weiterentwicklung ist: (1)~einen domänenspezifischen Detektor für Spieler, Torwart, Schiedsrichter und Ball trainieren, (2)~Teamzuordnung mit einer \emph{unknown}-Klasse und zeitlicher Glättung integrieren, (3)~Spielfeldkalibrierung samt Umgang mit Kamerabewegung implementieren, (4)~Ereignistabellen befüllen und (5)~Statistiken gegen manuell annotierte Referenzwerte validieren. Pixelgeschwindigkeiten dürfen dabei nicht als Meter pro Sekunde interpretiert werden; erst eine belastbare geometrische Abbildung erlaubt physikalische Größen.

\section{Fazit und Ausblick}

Der Prototyp demonstriert eine vollständige lokale Verarbeitungskette von der Videoquelle bis zur wiederverwendbaren synthetischen Personen-ID. Seine wesentliche Stärke liegt in der modularen Trennung von Detektion, Tracking, ReID, Qualitätsbewertung, Bewegung und Persistenz. Die Multi-Frame-Fusion und qualitätsgewichtete Profilaktualisierung adressieren ein praktisch relevantes Problem: ungeeignete Erstbeobachtungen sollen das Personenprofil nicht dominieren.

Der erreichte Stand ist als Applied-Sciences-Prototyp geeignet, um Entwurfsentscheidungen experimentell zu untersuchen. Eine belastbare Aussage über Erkennungsqualität oder Einsatzfähigkeit ist ohne annotierten Benchmark jedoch nicht möglich. Die nächsten priorisierten Schritte sind daher automatisierte Tests, ein versionierter Evaluationsdatensatz, reproduzierbare Abhängigkeiten, Profilversionierung und eine systematische Schwellwertkalibrierung. Parallel sind Datenschutzmaßnahmen technisch zu verankern. Erst danach sollten Skalierung über eine spezialisierte Vektordatenbank oder domänenspezifische Erweiterungen wie die Fußballanalyse erfolgen.

\begin{credits}
\subsubsection{\ackname}
\todoentry{Betreuung, Labor, Projektpartner und gegebenenfalls Förderhinweise ergänzen.}

\subsubsection{\discintname}
Die Autorin beziehungsweise der Autor erklärt, dass keine für den Inhalt dieser Arbeit relevanten Interessenkonflikte bestehen. \todoentry{Vor Abgabe prüfen und gegebenenfalls anpassen.}
\end{credits}

\begin{thebibliography}{13}

\bibitem{aharon2022botsort}
Aharon, N., Orfaig, R., Bobrovsky, B.-Z.: BoT-SORT: Robust Associations Multi-Pedestrian Tracking. arXiv:2206.14651 (2022). \doi{10.48550/arXiv.2206.14651}

\bibitem{dendorfer2021motchallenge}
Dendorfer, P., O\v{s}ep, A., Milan, A., Schindler, K., Cremers, D., Reid, I., Roth, S., Leal-Taix\'e, L.: MOTChallenge: A Benchmark for Single-Camera Multiple Target Tracking. Int. J. Comput. Vis. \textbf{129}, 845--881 (2021). \doi{10.1007/s11263-020-01393-0}

\bibitem{edpb-video}
European Data Protection Board: Guidelines 3/2019 on Processing of Personal Data through Video Devices, Version 2.0 (2020), \url{https://www.edpb.europa.eu/documents/guideline/guidelines-32019-on-processing-of-personal-data-through-video-devices_en}, last accessed 2026/09/02

\bibitem{gdpr}
European Parliament, Council of the European Union: Regulation (EU) 2016/679 (General Data Protection Regulation). Official Journal of the European Union L~119 (2016), \url{https://eur-lex.europa.eu/eli/reg/2016/679/oj}, last accessed 2026/09/02

\bibitem{luiten2021hota}
Luiten, J., O\v{s}ep, A., Dendorfer, P., Torr, P.H.S., Geiger, A., Leal-Taix\'e, L., Leibe, B.: HOTA: A Higher Order Metric for Evaluating Multi-object Tracking. Int. J. Comput. Vis. \textbf{129}, 548--578 (2021). \doi{10.1007/s11263-020-01375-2}

\bibitem{trackeval}
Luiten, J., Hoffhues, A.: TrackEval (2020), \url{https://github.com/JonathonLuiten/TrackEval}, last accessed 2026/09/12

\bibitem{mot17-data}
MOTChallenge: MOT17 -- Datensatz und Referenzannotation, \url{https://motchallenge.net/data/MOT17/}, last accessed 2026/09/12

\bibitem{ristani2016identity}
Ristani, E., Solera, F., Zou, R.S., Cucchiara, R., Tomasi, C.: Performance Measures and a Data Set for Multi-Target, Multi-Camera Tracking. arXiv:1609.01775 (2016). \doi{10.48550/arXiv.1609.01775}

\bibitem{streamlit-architecture}
Streamlit: Understanding Streamlit's Client-Server Architecture, \url{https://docs.streamlit.io/develop/concepts/architecture/architecture}, last accessed 2026/09/02

\bibitem{ultralytics-track}
Ultralytics: Multi-Object Tracking with Ultralytics YOLO, \url{https://docs.ultralytics.com/modes/track/}, last accessed 2026/09/02

\bibitem{zhang2022bytetrack}
Zhang, Y., Sun, P., Jiang, Y., Yu, D., Weng, F., Yuan, Z., Luo, P., Liu, W., Wang, X.: ByteTrack: Multi-Object Tracking by Associating Every Detection Box. In: Avidan, S., Brostow, G., Ciss\'e, M., Farinella, G.M., Hassner, T. (eds.) ECCV 2022. LNCS, vol.~13682, pp.~1--21. Springer, Cham (2022). \doi{10.1007/978-3-031-20047-2_1}

\bibitem{zheng2015market1501}
Zheng, L., Shen, L., Tian, L., Wang, S., Wang, J., Tian, Q.: Scalable Person Re-Identification: A Benchmark. In: Proc. IEEE Int. Conf. Comput. Vis., pp.~1116--1124 (2015). \doi{10.1109/ICCV.2015.133}

\bibitem{zhou2019osnet}
Zhou, K., Yang, Y., Cavallaro, A., Xiang, T.: Omni-Scale Feature Learning for Person Re-Identification. In: Proc. IEEE/CVF Int. Conf. Comput. Vis., pp.~3702--3712 (2019). \doi{10.1109/ICCV.2019.00380}

\end{thebibliography}

\end{document}
